%% kapa: exam=1 round=26 year=2015 period=1 form=A subject=civil_law subject_ko=민법 seq=12 number=12
%% source: Q-Net 감정평가사 기출문제 (cst00309) artlSeq=5138001 file=extracted/1008/5138001/제26회_감정평가사 1차_기출/제26회_감정평가사 1차_기출/1교시_감평_A형.pdf
%% conversion: kordoc 4.24.51 -> tools/gen1.py (auto); review=proofread-sonnet; flags=none
\begin{question}{12}
형성권에 관한 설명으로 옳지 않은 것은? (다툼이 있으면 판례에 따름)
\begin{choices}
  \item 형성권의 효력 발생에는 상대방의 동의나 승낙을 요하지 않는다.
  \item 형성권의 행사는 단독행위이므로 조건은 붙일 수 없음이 원칙이나, 계약의 정지조건부 해제는 인정된다.
  \item 공유물분할청구권은 형성권이다.
  \item 형성권은 반드시 재판상 행사해야 한다.
  \item 취소할 수 있는 법률행위의 취소권의 존속기간은 제척기간이다.
\end{choices}
\end{question}
\answer{4}
